%% ============================================================
%%  Soluzioni alle Domande di Revisione
%%  Dispensa del Corso di Patologia Post-Laurea
%%  Questo file è incluso con \input{} da main.tex; NON aggiungere il preambolo qui.
%% ============================================================

\chapter*{Soluzioni alle Domande di Revisione}
\addcontentsline{toc}{chapter}{Soluzioni alle Domande di Revisione}
\label{ch:solutions}

Le seguenti soluzioni forniscono la risposta corretta e una breve spiegazione per ciascuna delle
domande a scelta multipla presentate alla fine di ogni capitolo. Le spiegazioni sono concepite
per rinforzare i concetti chiave e chiarire i motivi per cui le risposte errate sono scorrette.

% ─────────────────────────────────────────────────────────────
\section*{Capitolo 1: Introduzione alla Patologia Digitale}
\addcontentsline{toc}{section}{Capitolo 1}

\begin{enumerate}

    \item \textbf{Risposta: C --- Immagini di vetrini interi (WSI) ad alta risoluzione.}\\
    La patologia digitale si basa sulla digitalizzazione dei vetrini istologici in immagini di
    vetrini interi (WSI) ad alta risoluzione, che possono essere archiviate, trasmesse e
    analizzate computazionalmente. Le opzioni A, B e D descrivono tecnologie o approcci non
    specifici della patologia digitale o non pertinenti alla sua definizione fondamentale.

    \item \textbf{Risposta: B --- Hamamatsu NanoZoomer, Leica Aperio, Philips IntelliSite.}\\
    Questi sono i principali produttori commerciali di scanner per vetrini interi (WSI) utilizzati
    nella patologia digitale. Le altre opzioni elencano strumenti di microscopia convenzionale,
    sistemi di imaging non correlati o combinazioni non corrette.

    \item \textbf{Risposta: C --- Circa 0,25 \textmu m per pixel.}\\
    Alla scansione a 40$\times$, la dimensione del pixel è tipicamente di circa 0,25~\textmu m,
    sufficiente a risolvere i dettagli subcellulari fini come la cromatina nucleare e le membrane
    cellulari. La scansione a 20$\times$ produce pixel di circa 0,5~\textmu m.

    \item \textbf{Risposta: B --- TIFF piramidale con tessere (ad es., SVS, NDPI, MRXS).}\\
    I formati WSI sono implementazioni proprietarie o aperte del TIFF piramidale con tessere, che
    consente la visualizzazione efficiente a diverse risoluzioni senza caricare l'intera immagine
    in memoria. I formati JPEG e PNG non supportano la struttura piramidale necessaria per le WSI.

    \item \textbf{Risposta: C --- OpenSlide.}\\
    OpenSlide è la libreria open source standard per la lettura di file WSI in formati proprietari
    (SVS, NDPI, MRXS, ecc.) in Python e altri linguaggi di programmazione. Le altre opzioni non
    sono librerie per la lettura di WSI.

    \item \textbf{Risposta: B --- Consente la visualizzazione efficiente a diverse risoluzioni
    senza caricare l'intera immagine in memoria.}\\
    La struttura piramidale delle WSI memorizza l'immagine a più livelli di risoluzione (livelli
    di zoom), consentendo ai visualizzatori di caricare solo le tessere necessarie per il livello
    di zoom corrente. Questo è essenziale per la gestione di file WSI che possono superare i
    100.000 $\times$ 100.000 pixel.

    \item \textbf{Risposta: C --- QuPath.}\\
    QuPath è la piattaforma open source più ampiamente utilizzata per l'analisi delle immagini di
    patologia digitale. Supporta la visualizzazione di WSI, l'annotazione, la segmentazione
    cellulare e tissutale, e l'integrazione con script Groovy e Python per l'analisi
    personalizzata.

    \item \textbf{Risposta: B --- Ematossilina ed eosina (E\&E).}\\
    La colorazione E\&E è il cardine della diagnosi istopatologica. L'ematossilina colora i nuclei
    di blu-viola e l'eosina colora il citoplasma e la matrice extracellulare di rosa. Tutte le
    altre colorazioni elencate sono colorazioni speciali o immunoistochimiche utilizzate per scopi
    diagnostici specifici.

    \item \textbf{Risposta: C --- Ridurre la variabilità cromatica inter-laboratorio per migliorare
    la generalizzabilità dei modelli di IA.}\\
    La normalizzazione della colorazione è una fase di pre-elaborazione che trasforma i colori
    delle WSI verso uno standard di riferimento comune, riducendo il \emph{domain shift} cromatico
    che può degradare le prestazioni dei modelli di IA quando applicati a dati di laboratori
    diversi da quelli di addestramento.

    \item \textbf{Risposta: B --- Apprendimento profondo, in particolare le reti neurali
    convoluzionali (CNN).}\\
    Le reti neurali convoluzionali (CNN) e le loro varianti (ResNet, EfficientNet, Vision
    Transformer) sono i metodi di apprendimento automatico dominanti nella patologia computazionale
    per compiti come la classificazione dei tumori, la segmentazione cellulare e la predizione
    degli esiti. I metodi tradizionali di apprendimento automatico (SVM, foreste casuali) sono
    stati in gran parte superati dall'apprendimento profondo per i compiti di analisi delle
    immagini.

\end{enumerate}

% ─────────────────────────────────────────────────────────────
\section*{Capitolo 2: Flusso di Lavoro Digitale e Controllo Qualità}
\addcontentsline{toc}{section}{Capitolo 2}

\begin{enumerate}

    \item \textbf{Risposta: C --- Formalina neutra tamponata al 10\%.}\\
    La formalina neutra tamponata al 10\% (FNT) è il fissativo standard per la processazione
    tissutale istologica di routine. Reticola le proteine preservando l'architettura cellulare e
    prevenendo l'autolisi. La glutaraldeide è utilizzata per la microscopia elettronica; l'alcool
    etilico è usato per alcuni campioni citologici; la soluzione di Bouin è un fissativo
    alternativo usato raramente nella pratica di routine moderna.

    \item \textbf{Risposta: B --- Ridurre la variazione nell'aspetto della colorazione causata da
    laboratori, protocolli e scanner diversi.}\\
    La normalizzazione della colorazione è una fase di pre-elaborazione computazionale che
    trasforma la distribuzione cromatica di un'immagine sorgente per farla corrispondere a quella
    di un'immagine di riferimento, riducendo la variabilità cromatica inter-laboratorio. Non
    aumenta la risoluzione, non elimina la necessità di IIC e non converte le immagini in scala
    di grigi.

    \item \textbf{Risposta: C --- Spessore tissutale non uniforme.}\\
    Lo spessore tissutale non uniforme --- causato da pieghe tissutali, bolle d'aria sotto il
    coprioggetto o un coprioggetto non uniforme --- impedisce all'algoritmo di messa a fuoco
    automatica dello scanner di mantenere la messa a fuoco sull'intera superficie del vetrino,
    producendo regioni sfocate nell'immagine. La sovra-colorazione e la contaminazione batterica
    non causano sfocatura; l'uso di un obiettivo a basso ingrandimento riduce la risoluzione ma
    non produce sfocatura nel senso tecnico.

    \item \textbf{Risposta: B --- Separare i contributi dei singoli coloranti (ad es., ematossilina
    e DAB) da un'immagine RGB.}\\
    La deconvoluzione cromatica di Ruifrok e Johnston utilizza la legge di Beer-Lambert per
    separare matematicamente i contributi dei singoli coloranti al colore di ogni pixel in
    un'immagine RGB. Questo consente l'analisi quantitativa separata dei canali dell'ematossilina
    (nucleare) e dell'eosina o DAB (citoplasmatico/membranoso).

    \item \textbf{Risposta: B --- Uno strumento open source per la valutazione automatizzata della
    qualità delle immagini di vetrini interi.}\\
    HistoQC è una pipeline open source e configurabile per il controllo qualità delle WSI,
    sviluppata da Janowczyk et al. Calcola metriche di qualità dell'immagine (messa a fuoco,
    intensità della colorazione, copertura tissutale, artefatti) e genera rapporti di qualità per
    vetrino. Non è un sistema commerciale, un protocollo di colorazione o un formato di file.

    \item \textbf{Risposta: C --- Tutte le fasi dalla ricezione del campione alla scansione del
    vetrino.}\\
    La fase pre-analitica comprende tutte le fasi che precedono la scansione: ricezione del
    campione, descrizione macroscopica, fissazione, processazione tissutale, inclusione,
    sezionamento, colorazione e coprioggetto. Queste fasi esercitano la maggiore influenza sulla
    qualità finale dell'immagine.

    \item \textbf{Risposta: B --- Il tempo di fissazione influenza la morfologia tissutale, la
    preservazione degli antigeni e le analisi IIC e molecolari successive.}\\
    Il tempo di fissazione è critico perché determina il grado di reticolazione proteica. La
    sotto-fissazione porta a scarso dettaglio nucleare e IIC inaffidabile; la sovra-fissazione
    causa una reticolazione eccessiva che può mascherare gli epitopi antigenici. Il tempo di
    fissazione non influenza la dimensione del file WSI, il numero di livelli della piramide o
    la velocità di scansione.

    \item \textbf{Risposta: C --- Riscansionare il vetrino dopo aver verificato la calibrazione
    della messa a fuoco dello scanner.}\\
    Il primo passo appropriato quando si scopre che una WSI presenta significative regioni sfocate
    è riscansionare il vetrino, dopo aver verificato la calibrazione della messa a fuoco dello
    scanner e ispezionato il vetrino per eventuali problemi fisici (pieghe, bolle d'aria). Accettare
    l'immagine sfocata per la diagnosi è inappropriato; eliminare l'immagine senza tentare la
    rimediazione è prematuro; applicare un filtro di nitidezza computazionale non recupera le
    informazioni perse a causa della sfocatura.

    \item \textbf{Risposta: B --- Garantire la tracciabilità dei campioni e dei vetrini attraverso
    il flusso di lavoro.}\\
    La documentazione della catena di custodia garantisce che ogni campione, cassetta, vetrino e
    immagine digitale possa essere tracciato attraverso il flusso di lavoro, dalla ricezione alla
    diagnosi. Questo è un requisito legale ed etico che protegge i pazienti dagli errori di
    identificazione del campione.

    \item \textbf{Risposta: B --- ISO 15189.}\\
    ISO 15189 è lo standard internazionale specifico per i laboratori medici, che specifica i
    requisiti per la qualità e la competenza nei test di laboratorio medico. ISO 9001 è uno
    standard generico per i sistemi di gestione della qualità; ISO 27001 riguarda la sicurezza
    delle informazioni; ISO 13485 riguarda i dispositivi medici.

\end{enumerate}

% ─────────────────────────────────────────────────────────────
\section*{Capitolo 3: Fondamenti dell'Intelligenza Artificiale nella Patologia}
\addcontentsline{toc}{section}{Capitolo 3}

\begin{enumerate}

    \item \textbf{Risposta: C --- Una rete neurale con molti strati nascosti che apprende
    rappresentazioni gerarchiche dei dati.}\\
    L'apprendimento profondo si riferisce alle reti neurali con molti strati nascosti (da qui
    ``profondo''), che apprendono rappresentazioni gerarchiche dei dati di input. Ogni strato
    apprende caratteristiche di complessità crescente, dalle caratteristiche di basso livello
    (bordi, texture) alle caratteristiche di alto livello (strutture cellulari, pattern tissutali).

    \item \textbf{Risposta: B --- Reti neurali convoluzionali (CNN).}\\
    Le CNN sono l'architettura di apprendimento profondo dominante per l'analisi delle immagini
    nella patologia computazionale. Utilizzano filtri convoluzionali per estrarre caratteristiche
    spaziali locali dalle immagini, rendendole particolarmente adatte all'analisi di immagini
    istologiche.

    \item \textbf{Risposta: C --- Apprendimento debolmente supervisionato con etichette a livello
    di vetrino.}\\
    L'apprendimento debolmente supervisionato (o apprendimento con supervisione debole) utilizza
    etichette a livello di vetrino (ad es., diagnosi del caso) anziché annotazioni pixel per pixel,
    riducendo il carico di annotazione. Questo è particolarmente rilevante nella patologia, dove
    le annotazioni dettagliate richiedono molto tempo da parte di patologi esperti.

    \item \textbf{Risposta: B --- Apprendimento per istanze multiple (MIL).}\\
    Il MIL è un paradigma di apprendimento debolmente supervisionato in cui un vetrino è
    rappresentato come un ``sacchetto'' di tessere (istanze), e l'etichetta del sacchetto (diagnosi
    del vetrino) è nota ma le etichette delle singole tessere non lo sono. Il modello apprende a
    identificare le tessere diagnosticamente rilevanti (istanze positive) all'interno del sacchetto.

    \item \textbf{Risposta: C --- Trasferimento dell'apprendimento da modelli pre-addestrati su
    grandi dataset di immagini naturali.}\\
    Il trasferimento dell'apprendimento (transfer learning) è la tecnica più comune per affrontare
    la scarsità di dati annotati nella patologia computazionale. I modelli pre-addestrati su
    ImageNet o altri grandi dataset vengono adattati (fine-tuned) su dataset di patologia più
    piccoli, sfruttando le rappresentazioni di caratteristiche apprese durante il pre-addestramento.

    \item \textbf{Risposta: B --- La capacità di un modello di generalizzare a nuovi dati non visti
    durante l'addestramento.}\\
    La generalizzazione si riferisce alla capacità di un modello di apprendimento automatico di
    applicare le conoscenze acquisite durante l'addestramento a nuovi dati non visti. Un modello
    con scarsa generalizzazione può avere prestazioni eccellenti sui dati di addestramento ma
    scarse sui dati di test, un fenomeno noto come overfitting.

    \item \textbf{Risposta: C --- Aumentazione dei dati.}\\
    L'aumentazione dei dati (data augmentation) è la tecnica di generare artificialmente nuovi
    esempi di addestramento applicando trasformazioni casuali (rotazione, riflessione, variazione
    del colore, ecc.) alle immagini esistenti. Questo aumenta la dimensione effettiva del dataset
    di addestramento e migliora la robustezza del modello.

    \item \textbf{Risposta: B --- Mappe di attivazione di classe (CAM) e metodi di attribuzione
    del gradiente (Grad-CAM).}\\
    Le CAM e i metodi basati sul gradiente come Grad-CAM sono le tecniche di interpretabilità più
    comunemente utilizzate nella patologia computazionale. Producono mappe di calore che
    evidenziano le regioni dell'immagine che hanno contribuito maggiormente alla decisione del
    modello, consentendo ai patologi di verificare se il modello si concentra sulle caratteristiche
    morfologiche corrette.

    \item \textbf{Risposta: C --- Bias nei dati di addestramento che riflette disparità nella
    raccolta dei dati o nelle pratiche di annotazione.}\\
    Il bias nei modelli di IA per la patologia deriva spesso da bias nei dati di addestramento,
    che possono riflettere disparità nella raccolta dei dati (ad es., sottorappresentazione di
    certi gruppi demografici), nelle pratiche di annotazione (ad es., variabilità inter-osservatore)
    o nelle pratiche di laboratorio (ad es., variabilità della colorazione).

    \item \textbf{Risposta: B --- Validazione esterna su dataset indipendenti di laboratori
    diversi.}\\
    La validazione esterna su dataset indipendenti di laboratori diversi è il gold standard per
    valutare la generalizzabilità di un modello di IA per la patologia. La validazione interna
    (su dati dello stesso laboratorio) può sovrastimare le prestazioni del modello a causa del
    \emph{domain shift} tra laboratori.

\end{enumerate}

% ─────────────────────────────────────────────────────────────
\section*{Capitolo 4: Apprendimento Profondo e Analisi Multimodale}
\addcontentsline{toc}{section}{Capitolo 4}

\begin{enumerate}

    \item \textbf{Risposta: C --- Vision Transformer (ViT).}\\
    I Vision Transformer (ViT) applicano il meccanismo di attenzione dei Transformer, originalmente
    sviluppato per il linguaggio naturale, alle immagini suddividendole in patch. Hanno dimostrato
    prestazioni competitive o superiori alle CNN su molti compiti di analisi delle immagini,
    inclusa la patologia computazionale.

    \item \textbf{Risposta: B --- Segmentazione semantica.}\\
    La segmentazione semantica assegna un'etichetta di classe a ogni pixel dell'immagine,
    consentendo la delineazione precisa di strutture tissutali (ad es., ghiandole, stroma, tumore)
    nelle WSI. Questo è distinto dalla classificazione (etichetta a livello di immagine) e dal
    rilevamento di oggetti (bounding box).

    \item \textbf{Risposta: C --- Integrazione di dati di imaging con dati genomici, clinici o
    di altra natura per migliorare la predizione degli esiti.}\\
    L'analisi multimodale nella patologia computazionale si riferisce all'integrazione di dati di
    imaging (WSI) con altri tipi di dati --- genomici, trascrittomici, clinici, radiologici ---
    per migliorare la predizione degli esiti rispetto a quanto possibile con un singolo tipo di
    dato.

    \item \textbf{Risposta: B --- Reti generative avversariali (GAN).}\\
    Le GAN sono il metodo di apprendimento profondo più comunemente utilizzato per la
    normalizzazione della colorazione e l'aumentazione dei dati nella patologia computazionale.
    Le CycleGAN, in particolare, consentono la traduzione di immagini non appaiata tra domini
    (ad es., da un laboratorio a un altro) senza richiedere coppie di immagini corrispondenti.

    \item \textbf{Risposta: C --- Predizione della sopravvivenza e degli esiti clinici da WSI.}\\
    La predizione della sopravvivenza e degli esiti clinici da WSI è uno dei compiti più
    clinicamente rilevanti nella patologia computazionale. I modelli di apprendimento profondo
    possono estrarre caratteristiche morfologiche prognostiche dalle WSI che non sono facilmente
    quantificabili dall'osservazione umana.

    \item \textbf{Risposta: B --- Modelli fondazionali pre-addestrati su grandi dataset di
    patologia (ad es., UNI, CONCH, Prov-GigaPath).}\\
    I modelli fondazionali per la patologia, pre-addestrati su milioni di tessere WSI mediante
    apprendimento auto-supervisionato, hanno dimostrato prestazioni superiori ai modelli
    pre-addestrati su ImageNet per molti compiti di patologia computazionale. Esempi includono
    UNI, CONCH e Prov-GigaPath.

    \item \textbf{Risposta: C --- Apprendimento auto-supervisionato su grandi dataset non
    etichettati.}\\
    L'apprendimento auto-supervisionato (self-supervised learning) consente ai modelli di
    apprendere rappresentazioni utili da grandi dataset non etichettati, senza richiedere
    annotazioni manuali costose. Questo è particolarmente vantaggioso nella patologia, dove le
    annotazioni richiedono molto tempo da parte di patologi esperti.

    \item \textbf{Risposta: B --- Attenzione multi-testa che modella le relazioni a lungo raggio
    tra le tessere WSI.}\\
    Il meccanismo di attenzione multi-testa dei Transformer consente di modellare le relazioni
    a lungo raggio tra le tessere WSI, catturando il contesto spaziale globale che le CNN con
    campi recettivi locali non possono facilmente catturare. Questo è particolarmente utile per
    compiti che richiedono la comprensione dell'architettura tissutale a scala multipla.

    \item \textbf{Risposta: C --- Fusione di caratteristiche estratte da WSI con dati genomici
    o clinici mediante reti neurali multimodali.}\\
    La fusione multimodale nella patologia computazionale combina caratteristiche estratte da
    diverse modalità di dati (WSI, genomica, clinica) mediante architetture di rete neurale
    progettate per integrare rappresentazioni eterogenee. Questo approccio ha dimostrato di
    migliorare la predizione degli esiti rispetto ai modelli unimodali.

    \item \textbf{Risposta: B --- Modelli linguistici di grandi dimensioni (LLM) e modelli
    visione-linguaggio (VLM) per la generazione di referti e il ragionamento diagnostico.}\\
    I modelli linguistici di grandi dimensioni (LLM) e i modelli visione-linguaggio (VLM) stanno
    emergendo come strumenti promettenti per la generazione automatica di referti patologici, il
    ragionamento diagnostico e l'interazione in linguaggio naturale con i sistemi di patologia
    digitale. Esempi includono GPT-4V, LLaVA e modelli specifici per la patologia come PathChat.

\end{enumerate}

% ─────────────────────────────────────────────────────────────
\section*{Capitolo 5: Etica, Bias e il Futuro della Patologia Digitale}
\addcontentsline{toc}{section}{Capitolo 5}

\begin{enumerate}

    \item \textbf{Risposta: C --- Bias algoritmico che porta a prestazioni diseguali tra gruppi
    demografici.}\\
    Il bias algoritmico nella patologia digitale si riferisce alla tendenza dei modelli di IA a
    avere prestazioni diseguali tra gruppi demografici (ad es., per razza, sesso, età), spesso
    a causa di una rappresentazione non equilibrata di questi gruppi nei dati di addestramento.
    Questo può portare a disparità nelle cure sanitarie se i modelli vengono implementati
    clinicamente senza un'adeguata valutazione dell'equità.

    \item \textbf{Risposta: B --- Interpretabilità e spiegabilità dei modelli di IA.}\\
    L'interpretabilità e la spiegabilità dei modelli di IA sono requisiti fondamentali per
    l'implementazione clinica sicura. I patologi devono essere in grado di comprendere e
    verificare le basi delle decisioni del modello per poterle integrare nella loro pratica
    diagnostica in modo responsabile.

    \item \textbf{Risposta: C --- Consenso informato, privacy dei dati e proprietà dei dati.}\\
    Le questioni etiche relative all'uso dei dati dei pazienti nella patologia digitale includono
    il consenso informato (i pazienti devono essere informati su come vengono utilizzati i loro
    dati), la privacy dei dati (protezione delle informazioni identificabili) e la proprietà dei
    dati (chi possiede i dati dei pazienti e chi può trarne beneficio economico).

    \item \textbf{Risposta: B --- Validazione prospettica in studi clinici controllati.}\\
    La validazione prospettica in studi clinici controllati è il gold standard per dimostrare
    l'utilità clinica di un sistema di IA per la patologia. La validazione retrospettiva su
    dataset storici, sebbene necessaria, non è sufficiente per dimostrare che il sistema migliora
    gli esiti clinici nella pratica reale.

    \item \textbf{Risposta: C --- Supervisione umana e responsabilità professionale del patologo.}\\
    Nonostante i progressi dell'IA, la supervisione umana e la responsabilità professionale del
    patologo rimangono fondamentali. I sistemi di IA sono strumenti di supporto decisionale, non
    sostituti del giudizio clinico del patologo. La responsabilità legale e professionale della
    diagnosi rimane con il patologo.

    \item \textbf{Risposta: B --- Regolamento sui Dispositivi Medici (MDR) e Regolamento sui
    Dispositivi Medico-Diagnostici In Vitro (IVDR).}\\
    Nell'Unione Europea, i sistemi di IA per la patologia digitale utilizzati per la diagnosi
    primaria sono regolamentati come dispositivi medici ai sensi del MDR (UE 2017/745) o come
    dispositivi medico-diagnostici in vitro ai sensi dell'IVDR (UE 2017/746), a seconda della
    loro funzione specifica.

    \item \textbf{Risposta: C --- Federated learning per l'addestramento di modelli su dati
    distribuiti senza condivisione dei dati grezzi.}\\
    Il federated learning consente l'addestramento di modelli di IA su dati distribuiti in più
    istituzioni senza condividere i dati grezzi dei pazienti, affrontando le preoccupazioni sulla
    privacy e sulla sovranità dei dati. I modelli vengono addestrati localmente in ogni istituzione
    e solo i parametri del modello (non i dati) vengono condivisi con un server centrale.

    \item \textbf{Risposta: B --- Integrazione dell'IA come strumento di supporto decisionale
    nella pratica diagnostica di routine.}\\
    Il ruolo più probabile dell'IA nella patologia nel prossimo futuro è come strumento di
    supporto decisionale che assiste il patologo nella diagnosi di routine, piuttosto che come
    sostituto del patologo. L'IA può aumentare l'efficienza, la riproducibilità e la sensibilità
    diagnostica, ma la responsabilità finale della diagnosi rimane con il patologo.

    \item \textbf{Risposta: C --- Disparità nell'accesso alle tecnologie di patologia digitale
    tra paesi ad alto e basso reddito.}\\
    Le disparità nell'accesso alle tecnologie di patologia digitale tra paesi ad alto e basso
    reddito sono una preoccupazione etica significativa. Sebbene la patologia digitale abbia il
    potenziale di migliorare l'accesso alla diagnosi esperta attraverso la telemedicina, i costi
    dell'infrastruttura tecnologica possono esacerbare le disparità esistenti se non vengono
    affrontate deliberatamente.

    \item \textbf{Risposta: B --- Collaborazione interdisciplinare tra patologi, informatici,
    bioinformatici e clinici.}\\
    Lo sviluppo e l'implementazione responsabile dell'IA nella patologia richiedono una
    collaborazione interdisciplinare tra patologi (che forniscono competenza di dominio e
    annotazioni), informatici e bioinformatici (che sviluppano e validano i modelli) e clinici
    (che definiscono le esigenze cliniche e valutano l'utilità). Nessuna singola disciplina può
    affrontare da sola le sfide tecniche, cliniche ed etiche della patologia digitale basata
    sull'IA.

\end{enumerate}
